\documentclass[10pt,a4paper]{amsart}
\usepackage[margin=19mm]{geometry}
\usepackage{amsmath,amssymb,mathtools}
\usepackage[scr=rsfs,cal=euler]{mathalfa}
\usepackage{xfrac}
\usepackage[unicode,hidelinks,bookmarks=false]{hyperref}
\newcommand{\ie}{i.e.\ }
\newcommand{\cf}{cf.\ }
\newcommand{\resp}{respectively\ }
\newcommand{\Subsection}{\subsection}
\providecommand{\nobreakdash}{\nobreak\hbox{-}}
\setlength{\emergencystretch}{2em}
\multlinegap0pt
\usepackage{amssymb}
\usepackage{mathtools}
\usepackage{float}
\usepackage{xcolor}
\usepackage{tikz}
\newtheorem{definition}{Definition}
\newtheorem{theorem}{Theorem}
\newcommand{\E}{\mathbb{E}}
\newcommand{\Pp}{\mathbb{P}}
\newcommand{\R}{\mathbb{R}}
\newcommand{\abs}[1]{\left\lvert #1\right\rvert}
\newcommand{\dd}{\,\mathrm{d}}
\newcommand{\sgn}{\operatorname{sgn}}
\newcommand{\spow}[2]{#1^{\langle #2\rangle}}
\title{Tail representations for absolute, signed-power,\\ inverse, and logarithmic moments}
\author{Roberto Vila, Eduardo Nakano, Cecilia Castro}
\date{Source: arXiv:2606.14019v2 (CC BY 4.0)}
\begin{document}
\maketitle

\noindent\emph{Source and licence.} Tail representations for absolute, signed-power,\\ inverse, and logarithmic moments. Version arXiv:2606.14019v2 (CC BY 4.0), distributed by arXiv under the Creative Commons Attribution 4.0 International licence, http://creativecommons.org/licenses/by/4.0/, as recorded in the arXiv licence metadata for this exact version and shown on the abstract page https://arxiv.org/abs/2606.14019v2. Full licence notice: \texttt{licenses/arxiv-2606.14019-CC-BY-4.0.txt}.\par\smallskip

\noindent\textit{Editorial scope.} Counted unit: the article's theorem thm:integer (source0.tex lines 422--447). The article's complete original proof of it is delivered with it (source0.tex lines 449--467, one \texttt{proof} environment of 19 lines). No mathematics is written, repaired or re-derived: the statement and the proof are the article's own lines, in the article's own notation, and no retained line is substituted or re-flowed.  Retained uncounted as the source-local context the proof needs: lines 164--185.  Nothing was omitted from any retained span, so the package carries no omission record.  No retained line contains a citation, so the wrapper carries no bibliography and no bibliography entry.  No retained line contains a figure environment, an image-inclusion command or drawing code, so no image is delivered and the package declares no asset dependency.  Editorial formatting only, in addition: an AMS \texttt{amsart} preamble that restates the article's own declarations and macros used by the retained lines, namely the environments \texttt{definition}, \texttt{theorem} on separate counters, together with the standard packages those lines need.  The retained environments print in this document as Definition 1, Theorem 1 in order of appearance, whereas the article prints the same items with the numbering of its own full document.  The complete native source of the article as distributed in the arXiv e-print for this exact version is bundled unchanged as \texttt{sources/202749/source0.tex}.
\begin{definition}\label{def:signed-power}
For $p>0$, the \emph{signed power} of $x\in\R$ is
\[
  \spow{x}{p}=\sgn(x)\abs{x}^p,
  \qquad \sgn(0)=0.
\]
The corresponding one-sided moments are
\[
  M_p^+(X)=\E[X_+^p],
  \qquad
  M_p^-(X)=\E[X_-^p],
\]
with values in $[0,\infty]$.  Define the absolute moment, possibly infinite,
by
\[
  A_p(X)=\E[\abs{X}^p]=M_p^+(X)+M_p^-(X).
\]
When $A_p(X)<\infty$, define the signed-power moment by
\[
  S_p(X)=\E[\spow{X}{p}]=M_p^+(X)-M_p^-(X).
\]
\end{definition}

\begin{theorem}[Integer-valued variables]\label{thm:integer}
If $X$ is integer-valued and $p>0$, then
\begin{align}
  M_p^+(X)
  &=\sum_{k=0}^\infty \Delta_p(k)\{1-F_X(k)\},
  \label{eq:integer-positive}\\
  M_p^-(X)
  &=\sum_{k=1}^\infty \Delta_p(k-1)F_X(-k).
  \label{eq:integer-negative}
\end{align}
Moreover,
\[
  A_p(X)=M_p^+(X)+M_p^-(X)
\]
as an identity in $[0,\infty]$.  If $A_p(X)<\infty$, then
\[
  S_p(X)=M_p^+(X)-M_p^-(X).
\]
In particular, for $p=1$,
\begin{equation}\label{eq:integer-mean}
  \E[X]
  =\sum_{k=0}^\infty\{1-F_X(k)\}
   -\sum_{k=1}^\infty F_X(-k),
\end{equation}
provided $\E[\abs X]<\infty$.
\end{theorem}

\begin{proof}
For an integer-valued variable, the relevant tail probabilities are constant
on every interval $(k,k+1)$, $k=0,1,\ldots$.  Hence
\begin{align*}
  M_p^+(X)
  &=\sum_{k=0}^\infty
    \int_k^{k+1}p t^{p-1}\Pp(X>k)\dd t\\
  &=\sum_{k=0}^\infty\Delta_p(k)\{1-F_X(k)\}.
\end{align*}
Similarly,
\begin{align*}
  M_p^-(X)
  &=\sum_{k=0}^\infty
    \Delta_p(k)\Pp(X<-k)\\
  &=\sum_{k=0}^\infty\Delta_p(k)F_X(-k-1).
\end{align*}
Reindexing the last series gives \eqref{eq:integer-negative}.  The remaining
claims follow from Definition~\ref{def:signed-power}.
\end{proof}

\end{document}
